% im Text
% \ac{<label>}		Returns full acronym, i.e. long (short), when first used, and only short form when used after first use.
% \acs{<label>}		Returns short form of acronym without parentheses.
% \acf{<label>}		Returns always full acronym, as in first use \ac.
% \acl{<label>}		Returns long form of acronym.
% plural variants:	append to any command above a "p" to access the respective plural form.
% \Ac{<label>}, \Acs{<label>}, \Acf{<label>}, \Acl{<label>}, or plural variants: Return the respective acronym form but making the first letter upper case

% \accite{<label>}		Returns reference labeled by the cite-key (APA-Style: without parentheses), see BIM-example in section 1.8.
% \acparencite{<label>}	for APA-Style: Returns reference with parentheses labeled by the cite-key, see BIM-example in section 1.8.
% \actextcite{<label>}	for APA-Style: Returns reference for use in text labeled by the cite-key, see BIM-example in section 1.8.

% ngerman options:  für Dativ und Genitiv-Formen kann die long form auch im entsprechenden Kasus ausgegeben werden, sofern long-<Kasus> definiert wurden, siehe GeoLab-Beispiel unten.
% 					\acd{<label>} bzw. \acg{<label>} gibt die dativ bzw. genitiv Version des Akronyms aus. Auch hier gibt es die Version mit Großbuchstaben.


% Use {} brackets if value contains any special characters or a komma!!
% You can also refer to another acronym/glossary entry:
\newabbreviation{2D}{2D}{zweidimensional}

\newabbreviation{3D}{3D}{dreidimensional}

\newabbreviation[
genitiv = Building Information Modelings,
dativ   = Building Information Modeling,
cite    = borrmann2021,  % kann mit \accite{<key>} im Text verwendet werden.
]{BIM}{BIM}{Building Information Modeling}

\newabbreviation{bsi}{bsi}{BuildingSMART International}

\newabbreviation{BGU}{BGU}{\DAfaculty}

\newabbreviation[genitiv = Computer Aided Designs]{CAD}{CAD}{Computer Aided Design}

\newabbreviation[
genitiv = {Geodätischen Labors},
dativ   = {Geodätischen Labor}
]{GeoLab}{GeoLab}{Geodätisches Labor}
